\documentclass[11pt]{article}
\usepackage[a4paper,margin=26mm]{geometry}
\usepackage{amsmath,amssymb,amsthm,booktabs,longtable,hyperref}
\newtheorem{theorem}{Theorem}
\newtheorem{lemma}[theorem]{Lemma}
\newtheorem{proposition}[theorem]{Proposition}
\newtheorem{corollary}[theorem]{Corollary}
\theoremstyle{remark}\newtheorem{remark}[theorem]{Remark}
\newcommand{\Q}{\mathbb Q}\newcommand{\Z}{\mathbb Z}
\newcommand{\Tr}{\operatorname{Tr}}\newcommand{\Nm}{\operatorname{Nm}}
\newcommand{\ord}{\operatorname{ord}}
\title{Audit and extension of the shared-product quartics}
\author{Z diophantine in Q project}
\date{19 November 2026}
\begin{document}\maketitle
\begin{abstract}
We audit the attached reductions of Gabe Conant with Astra~6.2 from 193 to 165 and from
229 to 180 rational coordinates, counting the distinguished input.
The exact supplied coefficient tables pass an independent expansion,
parameter-scaling, fixed-scalar and canonical-text verification.
We give the arithmetic transfer, with its hypotheses and boundary
cases explicit, and generalize the shared gate to any number of contact
arguments and any number field. The latter extension is a gate theorem,
not a new assertion that the necessary arithmetic inputs exist in every
number field. The same counts also support stronger fourth-order Lucas extraction,
which lowers the fixed 180-variable coefficient height from 161 to
108 bits. No smaller quartic variable count is claimed here.
\end{abstract}

\paragraph{Source credit.}
The supplied constructions and notes audited here are by Gabe Conant with
Astra~6.2, in the research bundles
\emph{H10-quartic-constructions} and
\emph{H10-eleven-variable-constructions}.
The audit, corrections, and further improvements stated in this paper are
subsequent project work. The original source files and their provenance
remain separately preserved.

\section{Status, sources and exact statements}
The sources are the user-supplied archive
\texttt{H10-quartic-constructions.zip}, in particular its
\texttt{quartics/} tables and transfer notes, and the project repository
\begin{center}\small
\url{https://github.com/JoseBrox/diophantine-z-q_improvements}
\end{center}
The pinned commit is
\begin{center}\small
\texttt{d3a8c223716a0fc0dbf5af4eb7998d640c21048d}.
\end{center}
We retain the exact arithmetic interfaces of the source 193-variable
manuscript v74 and, for the effective member, the 229-variable
manuscripts v62 and v68. These are hypotheses of the transfer below.
They are not conclusions of a finite coefficient computation.

The retained interfaces are as follows.
\begin{enumerate}
\item For $K=\Q(\sqrt2)$, $d=3-\sqrt2$, $E:y^2=x^3-d^2x$,
$R_0=(1+\sqrt2,2\sqrt2)$ and $P_0=2R_0$, the source group
and formal-logarithm statements hold. The charts give
$E(K)=\Z R_0+E[2](K)$ with the indicated label conventions.
The relative formal error has order four at every consumed good prime.
\item The scalar extraction and ordered scalar preparation hold.
For the general source $C=149^N$ and
$R_s=149^{2N+1}r/(1+149^2r)$, where
$\beta=51X^2+2504$ and
$r=\beta(\beta-1)(\beta-2401)(32\beta+343)$.
They exclude the exceptional reference strata before a ratio is used.
Their lower bound and scalar-height bound are the v74 bounds with
$L_N=\lfloor9(2N-1)/4\rfloor+1$.
\item The full-domain contact-height estimates hold for the prescribed
reflected ten-contact set, including the midpoint where the radical is
one. The general reflection has center $1$; the effective reflection
has center $3/22$. We do not substitute an arbitrary contact set.
\item The positive aggregate trace-height theorem holds on all the
finite point differences allowed by the charts, including finite
two-torsion differences.
\item The source pole obstruction excludes negative odd valuations
at the consumed primes. The prepared pole-realization theorem applies
after its actual rational argument has the required positivity, real
upper bound, square reduced denominator and local preparations at
$2,3,17$.
\item The fixed-data Gaussian prime-pattern and compatible annihilation
supplier permits arbitrary finite avoidance, arbitrarily large prime
cutoffs and arbitrarily large prescribed slope. Its real and finite
approximation statements are those used in the source ordered proof.
\end{enumerate}
The explicit theorem additionally uses the effective source
contact-height inequalities outside its stated exceptional sets, with
the numerical constants of v68. No new effective arithmetic estimate
is being inferred from the attached coefficient tables.

\begin{theorem}[Audited transfer]
Under the corresponding retained interfaces, the supplied balanced
165-variable quartic has rational input projection $\Z$ for every
sufficiently large fixed $N$. Its sufficient threshold remains
ineffective. For every integer $N\ge1$ it has degree exactly four,
139 quadratic rows, 2608 coordinate monomials, content one, and exact
height
\[
 \max\{126616122944180358034867798936,149^{2N+2}\}.
\]
Under the retained effective interfaces the supplied fixed
180-variable quartic has rational input projection $\Z$. It has degree
exactly four, 152 quadratic rows, 2648 monomials and height
\[
2584403986888404839167627077886434458396171597584.
\]
Its canonical literal has 54695 ASCII bytes and SHA-256
\begin{center}\small
\texttt{e6fd2cab3564dd0eb84c388aa252db1b6057f1135fbedfdf248416cd6585fd73}.
\end{center}
All counts include the distinguished input.
\end{theorem}

\section{The shared scalar gate, including its zero boundaries}
Let $m\ge2$ and $q_0,\ldots,q_{m-1}\in\Q^\times$. Put
$c=\sum_jq_j$ and $s_j=c-q_j$. The compiled chart requires $c\ne0$.
Let $W\in K$ be the common product of the field factors. Define the
rational $D$ by
\begin{equation}\label{eq:gate}
 D\Tr(W)=\Nm(W-1),\qquad
 g_j=1+D^2+s_j^2/q_j.
\end{equation}
The relation itself forces $\Tr(W)\ne0$: otherwise the norm vanishes,
so $W=1$ in the field $K$, contradicting $\Tr(W)=0$. Thus $D$ is
uniquely defined. This argument does not require a real sign
assumption on $W$.

At a consumed prime $p$ for argument $j$, write
\[
 v_p(q_j)=-\mu,\qquad v_p(R_j)=\mu,\qquad \mu>0\text{ odd}.
\]
The actual pole argument is
$b_j=(408/409)^2R_jg_j^2/4$.
The exceptional sets include $2,3,17$, so
$v_p(408/409)\le0$. This inequality also covers $p=409$;
the negative valuation there strengthens the obstruction.

\begin{lemma}\label{lem:gate}
Every rational solution satisfies
\[
 v_p(D)>-\mu/4,\qquad v_p(s_j)>-3\mu/4,
 \qquad v_p(c)=-\mu.
\]
If all field factors have nonpositive valuation at each slot above
$p$, then some slot satisfies $v(W)>-\mu/4$.
\end{lemma}
\begin{proof}
Set $v(0)=+\infty$. If $g_j\ne0$, the valuation of $b_j$ is odd,
and a negative odd valuation is forbidden. It follows that
$v_p(g_j)>-\mu/2$. The inequality also holds when $g_j=0$.
A negative finite valuation of $D^2$ is even; a finite valuation of
$s_j^2/q_j$ is $2v_p(s_j)+\mu$, which is odd. A negative minimum
therefore cannot cancel with the other term or the constant. Each
nonconstant term has valuation greater than $-\mu/2$, proving the
first two inequalities. Since $v_p(s_j)>-\mu$, the last follows.

Use the two conjugate valuation slots, with repetition when $p$ is
inert. If one slot has $v(W)=0$, it suffices. Otherwise both values
$w_1,w_2$ are negative. The shifted norm has valuation $w_1+w_2$,
whereas $v_p(\Tr W)\ge\min(w_1,w_2)$. Hence
\[
 v_p(D)\le\max(w_1,w_2),
\]
which proves the assertion. Zero components are omitted from minima;
$W=0$ is already impossible under \eqref{eq:gate}.
\end{proof}
No disjointness of contact radicals on an arbitrary rational
solution has entered this proof. Disjointness is used only in the
construction of integer witnesses.

\section{The four factors and all sixteen tests}
Let $A_i,B_i$ denote the source reciprocal abscissas, $D_H,B_0$ the
reference tests, $E_i$ the four cleared comparisons, $\mathcal H$ the
homogeneous Lucas residual, and $\mathcal N$ the cleared input
numerator. The four factors have constant term one and the following
nonconstant terms:
\[
\begin{array}{c|rrrr}
 &c&c^2&c^3&c^4\\\hline
0&\mathcal N^4&B_1^2&-E_1^4&D_H^4\\
1&E_4^4&A_3^2&-\mathcal H^4&B_2^4\\
2&E_3^4&B_4^2&-E_2^4&A_2^4\\
3&A_1^2&B_3^2&-A_4^2&B_0^4
\end{array}
\]
Reciprocal-abscissa valuations are even at good unramified slots.
For $v(c)=-\mu$ with $\mu$ odd, the first three displayed terms
belong to the distinct nonzero classes
$-\mu,-2\mu,-3\mu\pmod4$. The last term has class $4\pmod8$.
It cannot have valuation zero, and it cannot tie any of the first
three terms. Thus the minimum in each factor is unique, its
valuation is nonpositive, and a factor of valuation zero reduces to
one. The signs do not alter any of these assertions.

At the slot supplied by Lemma~\ref{lem:gate}, every factor and then
every term has valuation greater than $-\mu/4$, in particular greater
than the source cutoff $-\mu/2$. Every reciprocal consequently has
strictly positive even valuation, at least two. Every cleared error
has positive valuation, at least one. In the source notation
\[
 L_i-\overline{L_i'}=\frac{k_iE_i}{s_i\overline{t_i}},
 \qquad \mathcal H=s_4^2\mathcal G,
 \qquad L_1-\beta=\mathcal N/s_1.
\]
The positive reciprocal tests imply $v(s_i),v(\overline{t_i})\le-1$,
so the comparison and Lucas depths are at least three, and the input
depth is at least two. These divisions are made only after the local
nonvanishing has been established. Integrality is transferred through
one of the unit reference labels $h,h-1$ before relative formal
precision is used as absolute precision.

Consequently the source radical divisibilities survive:
\[
 M_j^3\mid\Delta,\qquad M_j^2\mid
 |\operatorname{num}(\beta-n)|
\]
when the latter numerator is nonzero. The source height proof requires
only the weaker first power in the second divisibility. If $M_j=1$
both divisibilities are automatic; no global quotient is imposed merely
to state them. Thus the allowed zero-ordinate and finite-torsion
boundaries are still covered by the source global height argument.

The coefficient savings use the exact signed identity
$c(U+cZ)(U-cZ)=cU^2-c^3Z^2$ for the first three rows.
The inverse map is $U=(P+Q)/2$, $Z=(P-Q)/(2c)$; the existing unit
condition makes this a rational bijection. In the last row one added
field coordinate $L=U+cZ$ gives
$cLU-c^2LZ=cU^2-c^3Z^2$.
These substitutions are applied in every retained physical equation,
including the different imaginary-component normalizations. They are
not identifications of independently chosen elliptic points.

\section{Integer completeness and the order of choices}
Fix an integer input. Choose the source scalar reference neighborhoods,
then the arbitrarily large positive scalar label $a\in6\Z$, the Lucas
labels, and the four nonzero nontorsion difference points. All their
multipliers, trace shifts and compact neighborhoods are now fixed.
For two arguments the shift is
$T=7+\Tr(\sum_i x(Q_i))/4>7$ on these witnesses. Indeed the first
two points are nonzero doubles, with $x>d$ at each real embedding,
and the other two have $x\ge-d$. The contact diameter is four.
For three arguments the successive shifts are at least seven and
nine and the contact diameter is $25/11$. Thus the translated contact
sets are disjoint before any comparison prime is chosen.

Increase the fixed avoidance cutoff to include poles and bad reductions
of the other weights evaluated at each translated contact. At a selected
prime for $j$, $v_p(q_j)=-1$ and every other weight is integral, giving
$v_p(c)=-1$.

The compatible Gaussian prime supplier may choose all selected primes
in a sufficiently short multiplicative interval. Each relevant reduced
elliptic group has order divisible by four and, by Hasse's bound, less
than twice the smallest selected prime. No selected odd prime divides
any such order: its only possible positive multiple below this bound
would be itself, incompatible with divisibility by four. Avoid also the
previous preparation moduli and fixed nonzero root multipliers.

The source synchronization now refines $h$ modulo powers of each
selected prime so that the primary and opposite reference formal depths
are equal. If the initial depths are $d_0,e_0$, only the smaller one is
raised, by imposing the appropriate valuation of $h$ or $h-1$.
One does not demand that both $h$ and $h-1$ have positive valuation.
The new prime-power constraints are coprime to inherited reduction
orders and to one another. CRT retains the annihilation conditions,
including the half-generator conditions. The source compact-component
argument retains the real returns on the refined progression.

At a successful slot of common depth $d\ge1$, reciprocals have depth
at least $2d$, cleared comparisons and the homogeneous Lucas error have
depth at least $2d$, and the input numerator has depth at least $3d$.
Every displayed nonconstant factor term therefore has positive
valuation. The shared $W$ is one modulo $p$ at that slot. At the
opposite slot every factor is either a unit of residue one or has a
pole. If $W$ is a unit there, $\Tr W\equiv2$ and the shifted norm
has positive valuation. If it is a pole of order $k$, the trace has
valuation $-k$ and the shifted norm has valuation at least $1-k$.
In either case $v_p(D)>0$, including $D=0$. Since $s_j$ is integral,
\[
 g_j\equiv1\pmod p,\qquad v_p(b_j)=1.
\]

At the prepared finite and real places, the reflected contact function
has the expansion
\[
 q(\lambda+u)=\kappa\lambda^{-2}+O(\lambda^{-3}),\qquad
 \kappa=\sum_{\theta\in\Theta_0}(\theta-a)^2>0,
\]
with the same leading coefficient for every fixed shift $u$ and
reflection midpoint $a$. Hence all $q_j\to0$, their pairwise ratios
tend to one, and $s_j^2/q_j\to0$. The bounded physical tests give
$W\to1$, $D\to0$, $g_j\to1$. This includes the place two; a bare
claim that numerator and denominator both vanish would not suffice.
The actual pole arguments consequently retain their prescribed local
square classes and tend to $(204/409)^2\in(0,17)$ at the real place.
The signed factors remain positive on sufficiently small prepared
neighborhoods by continuity. Positivity is not asserted for arbitrary
rational solutions.

The source contact factorization gives even valuation of each $R_j$
away from its own selected support. Multiplication by the rational square
$g_j^2/4$ preserves parity. The selected valuations are positive, so the
reduced denominators are squares. At this point the actual rational
arguments are fixed and meet the source realization hypotheses. Give
each pole its own witnesses, fill the complete finite charts, and apply
the invertible Laurent substitutions. No previously fixed prime or
root choice is revisited.

\section{Effective three-contact transfer and the numerical anchor}
The 41 retained scalar, reference and inverse-Vieta component rows match
the pinned source coefficient by coefficient. The remaining three pairs
use the v74 complete finite translated-point chart. Their difference
abscissas are $h_i(h_i+s_i+t_i)$, even when the difference is finite
two-torsion. The retained scalar primary ratio and scalar quotient force
the two denominators needed by scalar extraction; deleting its alternate
ratio permits a zero alternate ordinate but does not invalidate those
two denominators.

Writing $\sum_i x(Q_i)=S_0+\sqrt2S_1$, the two aggregate rows recover
exactly the effective source shifts
\[
 t_1-t_0=13+S_0/2,\qquad t_2-t_1=17+(S_0-S_1)/2.
\]
Thus the effective four-point estimate
$B_0^2<40\sum_{j<3}h(t_j)+2425$ and its label conventions are
unchanged. The source formal estimate at every prime outside $\Gamma$
recovers the required divisibilities, including $p=409$ as already
explained. With the original scalar, the source $N=15$ contradiction
therefore transfers unchanged.

The displayed fixed member instead uses
\[
 R_s=Ar/(1+25r),
\]
where
\[
\begin{split}
 A={}&23\cdot149\cdot263\cdot337\cdot631\cdot1063\cdot1093\cdot1171\\
 &{}\cdot1439\cdot2161\cdot2437\cdot2837\cdot2851\cdot3187
 \cdot3343\cdot3793\\
 ={}&645904574641049503754777843805872291231896003979.
\end{split}
\]
For each of these primes, exact exhaustive residue enumeration verifies
$r(X)(1+25r(X))\ne0$ for every $X\in\mathbb F_p$. If $v_p(X)<0$,
the leading coefficient $32\cdot51^4$ is a unit and
$v_p(r)=8v_p(X)=v_p(1+25r)$. Thus $v_p(R_s)=1$ for every rational
input, not only every integral input.

The forced prime 149 still excludes the exceptional reference strata
first. The relevant nonzero norms are $41^2/64$, $17^2/64$, $63^2/64$
and $7$, all 149-adic units. The unchanged scalar extraction then gives
$v_p(a)>9/4$ at each prime dividing $A$, hence $A^3\mid a$.
None of these primes belongs to the variable scalar radical $M_S(\beta)$.
At its primes the denominator $1+25r$ is a unit and the same extraction
continues to give $M_S(\beta)^3\mid a$. Therefore
\[
 A^3M_S(\beta)^3\mid |a|,
 \qquad A>149^{22},
\]
which is exactly strong enough for both effective source estimates
\[
 \log|a|\ge66\log149,\qquad
 h(\beta)<2^{238}149^{-66}|a|.
\]
The subsequent source height calculation uses only these bounds, the
retained shifts and the already proved divisibilities.

For integer inputs $r\ge1$, so $R_s\ge A/26>2^{20}$, the uniform
real scalar preparation bound. The numerator $Ar$ is coprime to
$1+25r$, so $R_s(R_s+9792)$ has square reduced denominator.
Uniform exact congruences are
\[
 r\bmod16\in\{4,8,12\},\qquad r\bmod3=1,
 \qquad r\bmod17=13.
\]
The denominator is a unit at these primes. Since
$9792=2^6\cdot3^2\cdot17$, the unit $1+9792/R_s$ belongs to
$1+8\Z_2$, $1+3\Z_3$, and $1+17\Z_{17}$, respectively.
The required local squares and the ordered scalar construction persist.

The coefficient table uses a split $A=UV$. Its parameter gauge fixes
input coordinate one and removes the parameter from all 151 non-forced
rows, up to nonzero row factors. In the forced row it changes
$1+22201x_{78}x_{90}-149C x_{78}$ to
$1+25x_{78}x_{90}-Vx_{78}$, with $C=U$.
This is exactly the physical replacement above. The remaining recorded
maps are input-fixing diagonal maps, invertible elementary row operations,
a witness permutation and integral shears. Every one is independently
checked. The final polynomial is the sum of the 152 resulting squares.

\section{Stronger extraction in the same numbers of variables}
The audit suggests a further improvement that is absent from the
tables of Gabe Conant with Astra~6.2. In the first three field factors, replace the two
signed terms $cU^2-c^3Z^2$ by
\[
                         c^5U^2-c^7Z^2.
\]
Keep the $c^2$ reciprocal square, the $c^4$ reciprocal fourth power,
and the entire fourth factor unchanged. The same hyperbolic variables
work, because
\[
              c^5(U+cZ)(U-cZ)=c^5U^2-c^7Z^2.
\]
The four residue classes used in the audit are unchanged. The shared
gate therefore still provides a slot with every term of valuation
greater than $-\mu/4$.

\begin{proposition}[One additional order of precision]
The modified factors give
\[
 M_j^4\mid\Delta,
 \qquad M_j^3\mid |\operatorname{num}(\beta-n)|
\]
when the corresponding residual is nonzero. They admit complete
quartic compilations in 165 and 180 variables, with 139 and 152
quadratic rows respectively. The balanced 165-variable family has
2607 monomials and the same exact coefficient-height formula as the
audited 165-variable family.
\end{proposition}
\begin{proof}
Each of the four cleared comparison errors and the homogeneous Lucas
error occurs to the fourth power with exponent five or seven in $c$.
Its valuation is thus greater than $19\mu/16$ or $27\mu/16$,
respectively. Since $\mu\ge1$, each error has valuation at least two.
The same holds for the input numerator, whose exponent is five.
The unchanged reciprocal tests still have even valuation at least two.
The comparison denominators and the squared homogeneous Lucas
factor regain two orders, while the input denominator regains one.
Thus the comparison and Lucas depths are at least four, and the input
depth is at least three.

As in the audit, transfer integrality before invoking the formal
logarithm. The retained relative error of order four is now an
absolute error of order four. Hence all four actual ratios are
congruent to their labels modulo $p^4$, the Lucas residual vanishes
modulo $p^4$, and the input residual vanishes modulo $p^3$.
Taking squarefree products proves the assertions. There is no
claim of greater label precision than the available order-four
formal estimate.

Completeness needs no additional formal depth. At the existing
synchronized depth $d\ge1$ the ordinary errors have valuation at
least $2d$ and the input numerator at least $3d$. The new error
terms have valuation at least $8d-7>0$; the new input term has
valuation at least $12d-5>0$. The opposite-slot separation, the
shared gate, and all fixed-place and real limits are unchanged.

The stored integral Laurent weights give ordinary quadratic rows
with exactly the same physical coordinates and the same three
unit coordinates $c,c^{-1},c^{-2}$. Their literal equations are
$c c^{-1}=1$ and $c^{-2}=(c^{-1})^2$. The independent checker
reverses both gauges and proves that exactly six factor-component
rows acquire the specified $c^4$ multiplier on their odd
hyperbolic term; every other physical row is unchanged. It also
expands the two systems and verifies the parameter balancing and
exact height for every $N\ge1$.
\end{proof}

\section{A smaller fully specified 180-variable member}
The stronger Lucas divisibility reduces the sufficient forced scalar
size substantially. Retain the effective source constants
\[
 \mathcal C<2^{300},\quad K=2^{304},\quad J_*<2^{75},
 \quad E(M)<(21/10)M^{1/6006}.
\]
With $B_0=\max(1,|n|,|a|,|t|,|b|)$ and $s=\log B_0$, the
same reflected-height calculation, now using $M_j^4\mid\Delta$,
gives, if $\Delta\ne0$,
\[
 B_0^2<121\mathcal C\sqrt{21/10}\,
 |\Delta|^{1/4+1/48048}
 \left(\log J_*+\frac18\log|\Delta|\right).
\]
The additive constants are absorbed exactly as in the effective
source: at $|\Delta|=1$ the same positive constant majorant applies.
Using $|\Delta|\le14B_0^4$ gives
\begin{equation}\label{eq:improved-endpoint}
 e^{(1-1/12012)s}<350\,2^{300}(53+s/2).
\end{equation}
For example, the exact rational estimates
\[
 \sqrt{21/10}<29/20,\quad 14^{1/4}<97/50,
 \quad 14^{1/48048}<10001/10000
\]
show the coefficient is less than 350. The last estimate follows
already from the first three terms of the binomial expansion of
$(1+1/10000)^{48048}$. Also $\log2<7/10$ and $\log14<4$ imply
the displayed bound on its linear factor.

The quotient of the left side of
\eqref{eq:improved-endpoint} by $53+s/2$ is increasing for $s\ge0$,
since $12011/12012>1/106$. At $s=316\log2$, division by $2^{300}$
and $e^{-x}>1-x$ give the exact reversed comparison
\[
 2^{16}\left(1-\frac{316\cdot7}{10\cdot12012}\right)
 =\frac{137986048}{2145}
 >57260
 =350\left(53+\frac{158\cdot7}{10}\right).
\]
Thus $B_0\ge2^{316}$ excludes $\Delta\ne0$.

Choose the smaller squarefree anchor
\[
\begin{split}
 A_\mathrm{new}={}&23\cdot149\cdot263\cdot631\cdot1063\cdot1171
 \cdot1439\cdot2161\cdot2437\cdot2851\cdot3343\\
 ={}&51132150162547763354407979038657.
\end{split}
\]
Every prime is one of the previously certified anchor primes, and
\[
 A_\mathrm{new}^3\ge2^{316},\qquad A_\mathrm{new}>26\,2^{20}.
\]
With $R_s=A_\mathrm{new}r/(1+25r)$, the same reference exclusion,
scalar extraction and ordered scalar completeness proof applies.
It gives $c_a=|a|\ge A_\mathrm{new}^3$ and
\[
 \tau=h(\beta)<2^{238}A_\mathrm{new}^{-3}c_a<c_a/2.
\]
Consequently $\Delta=0$ by the preceding endpoint.

For the remaining input residual it is enough to use the weaker
source divisibility $M_j\mid J_X$, which follows from the improved
third power. The unchanged coarse calculation gives
\[
 s<\frac{23}{22}\log(256036K)+\frac{12}{11}\tau<240+c_a.
\]
Here the constant inequality follows from
$(23/22)\,322\,(7/10)<240$, while the coefficient of $c_a$ is
less than one. The source Lucas descent gives
$s\ge(c_a-6)\log(c_a-3)$. Since $c_a\ge2^{316}$, this is
strictly greater than $2(c_a-6)>240+c_a$, a contradiction.
Therefore $J_X=0$, and the source squarefree-denominator argument
again gives $X\in\Z$. Integer completeness was proved above.

\begin{theorem}[Improved fixed member]
Under the same effective arithmetic interfaces, the improved literal
specifies an integer quartic in 180 total rational variables whose
input projection is $\Z$. It has 152 quadratic rows, 2850 monomials,
content one, and exact coefficient height
\[
             206625078319751697390503348692624,
\]
which has 108 bits, compared with the supplied member's 161 bits.
Its canonical literal has 59084 ASCII bytes and SHA-256
\begin{center}\small
\texttt{5ee7330c9e2edd7d0f918fc692e7824089007ee0358e47079601e8aa9c8b4683}.
\end{center}
\end{theorem}
\begin{proof}
Use the input-fixing parameter gauge of the checked enhanced table
and the split
\[
 U=3593614808933267,\qquad V=14228611824350171,
 \qquad UV=A_\mathrm{new}.
\]
The forced row in this gauge has constant monomial $c^2$ rather
than 1; replacing its two scalar coefficients still gives exactly
$R_s=A_\mathrm{new}r/(1+25r)$. All other rows have a common parameter
power under the stored physical gauge, so the substitution changes
no other physical equation. Apply the recorded diagonal parameter
weights and a witness permutation. Finally multiply each of the
three contact-combination coordinates by 245 and primitively
normalize the affected rows. This is an invertible rational diagonal
map and fixes the input. The independent checker reconstructs all
these operations from the enhanced source, expands every row
square, and verifies the exact maximum coefficient, equal to $16U^2$,
and the literal. The input fourth-power coefficient remains one.
\end{proof}
The new member prioritizes smaller coefficient height. The supplied
180-variable member remains a valid alternative with fewer monomials
and a slightly shorter literal.


\section{A general number-field gate theorem}
The quadratic-field mechanism admits a useful precise extension.
For a number field $F/\Q$ of degree $d\ge2$, let
$e_{d-1}(W)$ denote the elementary symmetric function of degree $d-1$
of the $d$ conjugates of $W$; it is rational and has polynomial degree
$d-1$ in field coordinates. Define
\begin{equation}\label{eq:general-gate}
 D\,e_{d-1}(W)=\Nm_{F/\Q}(W-1).
\end{equation}
For $d=2$, this is exactly the preceding trace gate.

\begin{theorem}[Shared gate over an arbitrary number field]
Suppose $p$ is unramified in $F$ and the pole obstruction and rational
contact valuations are as in Lemma~\ref{lem:gate}. Assume all factors
of $W$ have nonpositive valuation at all $d$ conjugate slots, counted
with multiplicity. Use \eqref{eq:general-gate} and the same
$g_j=1+D^2+(c-q_j)^2/q_j$. Then its denominator is automatically
nonzero on every rational solution, and some slot satisfies
$v(W)>-\mu/4$.

For completeness, suppose the prepared $W$ at every slot is either a
unit congruent to one or has a pole, with at least one unit slot.
If $p>d$, then $v_p(D)>0$. Thus the same integer-witness preparation
works for any fixed number $m$ of separated contact arguments, provided
its remaining source interfaces and common weight asymptotic hold.
\end{theorem}
\begin{proof}
If $e_{d-1}(W)=0$, the equation forces $W=1$ because the norm is the
norm of a field element. But $e_{d-1}(1)=d\ne0$, a contradiction.
The rational parity proof of Lemma~\ref{lem:gate} is unchanged.
If one slot has nonnegative valuation, use it. Otherwise all slot
values $w_i$ are negative. Every summand of $e_{d-1}(W)$ has valuation
at least $\sum_iw_i-\max_iw_i$. The shifted norm has valuation
$\sum_iw_i$. Hence $v_p(D)\le\max_iw_i$ as required.

For the preparation, $W\ne0$ and
$e_{d-1}(W)=\Nm(W)\Tr(W^{-1})$.
If precisely $k$ slots are units congruent to one, then
$\Tr(W^{-1})\equiv k\pmod p$. Here $1\le k\le d<p$, so the trace
is a unit. At a pole slot $v(W-1)=v(W)$, and at a unit slot
$v(W-1)>v(W)=0$. Consequently the numerator in
\eqref{eq:general-gate} has strictly greater valuation than its
denominator, giving $v_p(D)>0$.
At fixed real and finite places $W\to1$ gives $D\to0$ because the
denominator tends to $d$. The common asymptotic gives
$(c-q_j)^2/q_j\to0$ for every fixed $m$. The remaining preparation
is precisely the already stated source interface.
\end{proof}
\begin{remark}
The field hypothesis is essential for automatic denominator exclusion:
in a product of fields a zero norm need not force $W=1$ in every
component. The extension does not supply elliptic, prime-selection,
contact-height or pole-realization inputs for new number fields. Its
gate has degree $d$ before any quadratic compilation, so a new
quartic count is not implied for $d>2$.
\end{remark}

\section{Verification limits and further searches}
The independent standard-library checker reads the supplied tables,
expands every row square, checks all active coordinate indices and
input fourth-power coefficients, derives the exact general height
for all $N\ge1$, applies every fixed-map operation afresh, and parses
the literal independently. It also checks 26778 residues in the 16
anchor prime fields and the complete modulo-eight noncancellation
classes for every odd contact depth. The supplied physical checker
separately reverses every Laurent scaling and compares to the pinned
physical source. Its scalar checker matches all 41 retained source
rows. These are complementary checks with explicitly different scopes.

Two documentation references need normalizing when integrated:
\texttt{scalar\_audit.py} in the transfer prose means the supplied
\texttt{check\_scalar.py}, and the example \texttt{check\_gated.py} in
the physical checker's docstring means \texttt{check\_physical.py}.
No mathematical repair of the two claimed reductions was needed in
this audit.

Exploratory linear-integer searches found no feasible unit weighting
that deletes the stored square of $c^{-1}$ from the supplied packing,
or that turns its last partial hyperbolic pair into a full pair with
only the same three unit coordinates. These observations are not
certificates of optimality and are not lower bounds on arbitrary
quartic presentations. A different chart, packing, or arithmetic
mechanism could lower the counts further.
\end{document}
